
% To further validate this result quantitatively, we selected 18 historically significant abstract (non-mainstream nowadays) artworks from artchive.com, captioned them, and evaluated them using each reward model. The selection of only 18 real artworks reflects a deliberate trade-off: we sacrificed scale to guarantee quality. Each work was chosen for its documented significance in art history and its clear deviation from conventional aesthetics — ensuring our evaluation targets meaningful counterexamples, not noise. We compared the value distribution using the Mann-Whitney U with an alternative hypothesis that these real arts will have lower scores. We calculated the $p$-value and $d$ effect size, as shown in Table~\ref{tab:real_art_scores}. We observed that there is both a statistical and a practical significant difference between the two distributions. This confirms the finding from LAPIS \cite{maerten_lapis_2025} that abstract arts usually have a lower rating than figurative art, and using reward models that only give a single value for each image, it might further propagate this bias. This also shows that these reward models do not actually understand what makes good art, but are merely a population opinion estimator. We ruled out the possibility that the model cannot understand abstract images by testing BLIP scores with images and captions from real art, and all of them have a $>0.99$ BLIP score, which shows that reward models could potentially understand abstract images, but chose to score them low because they diverged from mainstream. 

% Besides reward models, we also tried to use image generation models to reproduce these anti-aesthetics artworks by using the prompt generated. In Figure~\ref{fig:real_arts_gen}, we compared side-by-side with the selected original art and two generated images from two models. We also selected Flux Krea and DanceFlux as examples. Although sometimes their performance is close, we can see in some cases that Flux Krea can reproduce images that are much closer to the intended style.

% \begin{table}[] 
%     \centering
% \small
% \begin{tabular}{l|rr}
%   Reward Model & p-value & d \\ \hline
%  HPSv3 & $<1\times 10^{-20}$& -0.91\\
%  HPSv2 & $<1\times10^{-20}$& -0.95\\
%  ImageReward & $<1\times 10^{-20}$& -0.69\\
%  BLIP& 1.0&0.63\\
% \end{tabular}
%     \caption{Statistical test (with an alternative hypothesis that real arts have a lower score compared to AI-generated images) between the score distribution of the original image with original prompts and real historically significant abstract arts scored by each model.}
%     \label{tab:real_art_scores}
% \end{table}



\subsection{Validation On Real Arts}
Although these image reward models are primarily designed to evaluate AI-generated imagery, we also tested them on real artworks. Although these reward models might be meant to assess AI images, we tested them on real art to see how they would rate such works if generated by AI in the same style and meaning. This can validate our hypothesis that reward models are biased against non-mainstream art. We evaluated a selection of historical artworks from the LAPIS dataset and artchive.com. This approach also allows us to determine if the low scores are specific to our generated anti-aesthetic images or if they reflect a broader bias. We conducted a small-scale quantitative analysis first to examine each artwork individually to ensure. The selected artworks and their assigned scores are presented in Figure~\ref{fig:real_art_scores}.

For each artwork, we used \texttt{Qwen/Qwen3-VL-235B-A22B-Instruct} to generate a factual caption, while excluding interpretive content. The image and its corresponding caption were then fed to each reward model for rating. To provide a baseline for these scores, Table~\ref{tab:reward_model_baselines} lists the mean and standard deviation of scores from each reward model using original prompts on the original images generated by models we tested.

 We can observe that some of these scores are lower than 2 standard deviations from the mean. This confirms our theory that these reward models are heavily tuned for a general human preference and overlook the values of non-mainstream aesthetic images. 

To further validate this result quantitatively, we selected about 10K real artworks from the LAPIS Dataset \cite{maerten_lapis_2025}, which covers many styles and genres. We captioned them using \texttt{Qwen/Qwen3-VL-30B-A3B-Instruct} and evaluated them using each reward model. We compared the scores of real artworks with the benchmark each model released, which are usually popular models at the time of release. We report the raw score ($r$), relative score ($s$) within the leaderboard range, the rank on the leaderboard ($R$), the model that scored right above these real arts ($M'$), and the month the benchmark is released. The relative score is calculated by $(r_{art}-r_{min})/(r_{max} - r_{min})$, where $r_{art}$ is the score real art images get, and $r_{max}$ and $r_{min}$ are the highest and lowest scores in the leaderboard. Results are in the Table~\ref{tab:real_arts}. We can observe that the real art ranked very low in all three reward model leaderboards, and even lower than many early-stage image generation models. Additionally, they are much lower than the original images we generated that were scored with the original prompts. This shows that the reward models are heavily biased against. 

This also shows that these reward models do not actually understand what makes good art, but are merely a \textit{sub}-population (of the labeller) opinion estimator. We ruled out the possibility that the model cannot understand abstract images by testing BLIP scores with images and captions from real art. We achieved a BLIP score of 0.996. To further rule out that BLIP assigns a high score regardless of the prompt, we shuffled the prompt and images, resulting in a BLIP score of 0.086. This BLIP confirmation shows that reward models could potentially understand abstract images, but chose to score them low because they diverged from mainstream. 
\begin{table}[]
    \centering
    \small
    \begin{tabular}{c|ccclc}
    \toprule
      RM & $r$ & $s$ & $R$ &  MM/YY &$M'$ \\ \midrule
      ImgRewd & -0.13 & 0.74 & 5/7&  23/04 &SD1.4\\
 HPSv2& 0.21& 0.37& 18/23&  23/06&DALL·E mini\\
 HPSv3& 5.86& 0.57& 10/12&  25/08&Hunyuan-DiT\\
 \bottomrule
    \end{tabular}
    \caption{How reward models rate real art images.}
    \label{tab:real_arts}
\end{table}
